% Image credits for the photographs and AI illustrations of Book 5.
% Machine-readable license records live in images/book5/CREDITS.md; keep the
% two files in sync. The Book 5 writer owns this file: add one \item per
% photograph (inside the itemize, by chapter) as photographs land, and remove
% the "(none yet)" item with the first one.
\chapter*{Image Credits}

The drawings, diagrams and charts in this book are original; charts drawn
from published data name their source in the caption or in this list.
Photographs are used under their respective free licenses, with thanks to
their authors; full source links and license URLs for every photograph are
listed in \texttt{images/book5/CREDITS.md} in the book's repository.

\bigskip
Alongside the photographs and the original drawings, some figures are
AI-generated illustrations, produced with OpenAI image generation from
prompts written by the book's editors, and reviewed for accuracy before
inclusion. They show generic scenes only, never a real person, a real
product or a real brand. The full prompt for every generated image is
recorded in \texttt{images/book5/ai/PROMPTS.md} in the repository.

\bigskip
\noindent Photographs, by chapter:
\begin{itemize}\small
  \item Chapter 1: an IBM 704 at the NACA Langley laboratory, 1957 (NASA photograph GPN-2000-001881, public domain).
  \item Chapter 9: a magnetic-core memory plane of 32 by 32 cores, by Konstantin Lanzet, 2009 (Wikimedia Commons, CC BY-SA 3.0).
  \item Chapter 12: an opened 2.5-inch hard disk, by Evan-Amos, 2013 (Wikimedia Commons, CC BY-SA 3.0).
  \item Chapter 15: the Pilot ACE in the Science Museum, London, by Philafrenzy, 2011 (Wikimedia Commons, CC BY 2.0).
  \item Chapter 24: a Cray-1 in the Deutsches Museum, Munich, by Clemens Pfeiffer, 2006 (Wikimedia Commons, CC BY 2.5).
  \item Chapter 27: an Enigma machine in the Museo della Scienza e della Tecnologia, Milan, by Alessandro Nassiri, 2012 (Wikimedia Commons, CC BY-SA 4.0).
\end{itemize}

\noindent The four rendered images of Chapter~31 are written by the chapter's own C programs, in the repository under \texttt{code/bachelor-3/31-computer-graphics}.
